\documentclass[instructions.tex]{subfiles}
\begin{document}
\chapter{On perd son âme par le péché mortel.}


Le péché mortel déshonore l’âme, il la défigure, il lui donne la mort.


\primo Oui, le pécheur par sa conduite se déshonore jusqu’à se rendre semblable aux bêtes insensées\footnote{Homo cum in honore esset, non intellexit : comparatus est jumentis insipientibus. Ps. 48. 13.}. Il n’est plus, à proprement parler, un homme, mais un homme tout animal, dit saint Paul : \textit{Animalis homo} ; parce qu’il suit en aveugle ses inclinations, comme les animaux\footnote{Similis factus est illis. Ibid.}, et use de sa raison plus mal que les bêtes n’usent de leur instinct. Le péché a transformé l’ange en démon, et le péché change l’homme et le réduit à l’état des brutes. Le démon a péché en voulant se rendre semblable à Dieu ; et l’homme, dont l’âme est l’image de Dieu, en péchant se rend semblable aux bêtes, en s’attachant aux plaisirs, aux objets sensibles, comme les bêtes. Âme chrétienne, oubliez-vous ce que vous êtes ? À quoi vous réduisez-vous ? Ô mon Dieu ! devrais-je dire avec plus de confusion que David, comment oserai-je paraître devant vous, m’étant réduit par le péché à l’état d’une brute insensée\footnote{Ut jumentum factus sum apud te. Ps. 72. 23.} ?


\secundo L’âme est donc déshonorée par le péché mortel ; et par une suite nécessaire, elle est tellement défigurée, dit saint Augustin, qu’elle est plus horrible et plus insupportable à Dieu que l’infection des tombeaux ne l’est à la personne la plus délicate. Vous ne voulez pas le comprendre à présent ; mais vous le comprendrez peut-être trop tard. Lorsqu’une âme est séparée de son corps, si elle a le malheur d’être en péché mortel, la vue de sa laideur et de sa propre difformité la jette dans une confusion si accablante, qu’elle-même se précipiterait en enfer plutôt que d’entrer dans le ciel et de paraître devant la sainteté de Dieu en cet état.


\tertio Par le péché mortel on défigure son âme, et on lui donne la mort. Toute âme, dit le Seigneur, qui aura péché mourra\footnote{Anima quæ peccaverit, ipsa morietur. Ezech. 18. 4. 20.}. \textit{Ses péchés}, dit l’Écriture, \textit{sont comme des lions, dont les dents meurtrières lui donnent la mort}. Il est vrai qu’elle a toujours sa vie et son être naturel, parce qu’elle n’est pas détruite ; mais elle n’en est pas moins morte, parce qu’elle n’a plus la vie surnaturelle. Un cadavre est toujours un corps, il ne cesse pas d’être ; mais c’est un corps qui est mort. De même une âme en péché mortel est toujours une âme, et ne cesse pas d’être ; mais c’est une âme morte. Vous paraissez vivant aux yeux des hommes ; mais, sans la grâce de Dieu, vous êtes aussi véritablement mort que les cadavres qui sont dans les tombeaux\footnote{Nomen habes quod vivas, et mortuus es. Apoc. 3.}.

Votre plus grand malheur serait d’être insensible sur un état si déplorable. Un homme de distinction, nommé Sabinien, ayant séduit une vierge, saint Jérôme en versa des larmes, et lui écrivit ces paroles : « Malheureux, loin de pleurer sur la noirceur de votre crime, vous n’en rougissez même pas ! Voilà que je pleure de ce que vous ne vous pleurez pas vous-même, de ce que vous ne sentez pas que vous êtes mort\footnote{Hoc plango quod te ipsum non plangis ; quod te non sentis mortuum. Ad Sabin.}. » Ô pécheur, ouvrez les yeux sur votre misère. Serez-vous toujours ennemi de vous-même ? Ayez enfin pitié de votre âme.


\section{Résolutions}

\primo Je m’humilierai profondément à la vue de la laideur dont j’ai couvert mon âme par mes péchés mortels.

\secundo Je m’efforcerai de lui rendre la beauté dont la grâce sanctifiante la décorait.

\tertio Les sacrements, la prière, la pénitence, les bonnes œuvres, l’application à mes devoirs, voilà les moyens que je prendrai pour l’orner aux yeux de Dieu.

Faites, ô mon Dieu ! que je purifie entièrement mon âme des souillures que le péché y a imprimées, afin qu’elle redevienne digne de vous.
\end{document}
